%\section{\textit{\toolname{}} in Detail}
%\label{sec:approach}

%This section provides details of \toolname{}. As mentioned above, \toolname{} takes as input a to-be-verified data-structure module \texttt{code} together with its accompanying unit test suite \texttt{test}. \toolname{} keeps the test suite unchanged, and adds annotations to the data-structure implementation \texttt{code}. Finally, \toolname{} outputs the verified version of data-structure module.


\section{Stage 1: Generating the Annotations}
\label{sec:genfunc}
\begin{small}
\begin{algorithm}[t]
\small
%\SetAlgoLined
\SetAlgoNoLine 
\SetKwProg{Fn}{Function}{}{end}
\caption{The Generation Stage}
\label{fig:genfunc}
\textbf{Parameter.} {\texttt{n}: Number of samples generated in each module invocation.}\\
\textbf{Input.} {\texttt{code}: the data structure code to be verified; \texttt{test}: unit test suite}\\
\textbf{Output.} {\texttt{code} with annotations generated}\\
\Fn{{{\rm \genfunc{}}}{\rm (}{\tt code}, {\tt test}{\rm )}}{
\texttt{plan} $\gets$ {\textbf{ExecPlanner}}(\texttt{code})\\
\ForEach{$M_i$ $\in$ {\tt plan}}{
    \texttt{code} $\gets$ {\textbf{ExecWithSampling}}($M_i$, {\tt code}, {\tt test}, {\tt n})
}
\Return {\tt code}\\
}
\end{algorithm}
\end{small}

%This section details the first stage---annotation generation. As outlined in \Cref{sec:overview}, this process involves producing four categories of annotations: ($M_1$) the \texttt{View} trait, ($M_2$) the type invariant, ($M_3$) the specifications (pre- and post-conditions), and ($M_4$) a proof blocks module for generating proof blocks and loop invariants. 
%\subsection{The Generation Stage}

The generation stage creates annotations. As mentioned above, we may need to generate different kinds of annotations, so \toolname{} uses four dedicated modules, one for each type of annotation: ($M_1$) a view module, ($M_2$) a {type invariant} module, ($M_3$) a specifications module (for specification functions, preconditions, and postconditions), and ($M_4$) a {proof blocks} module, which generates both proof blocks and invariants.

%Each category $M_i$ ($i \in {1, \dots, 4}$) is generated by a dedicated module $G_i$, responsible for producing all annotations within its respective category.%\toolname{} provides a dedicated generation module for each generating each of these. %Calling a module will generate all annotations corresponding to the invoked module.

By examining the dependencies among these kinds of annotations, we observe that it should typically be sufficient to invoke the modules in order: $\langle M_1, M_2, M_3, M_4\rangle$. At the same time, it is not always necessary to execute every module. For example, some data structures do not require non-trivial type invariants; in such cases, generating an invariant may incur: (1) more instability in the language model (and thus affect the robustness of our framework); (2) unnecessary computational cost, and (3) more runtime overhead. To mitigate these issues, we employ a \emph{planner} module that determines which generation modules should be executed.

\paragraph{Planner} The generation process begins with the planner agent (Line 5). The function \textbf{ExecPlanner} accepts the target code as input, invokes the planner agent, and outputs the modules to execute.
The planner agent instructs the LLM to select only the necessary generation modules to be executed. Its prompt comprises four blocks: (i) a concise description of the role of the planner; (ii) a catalog of the available generation modules: there are four in our current framework; (iii) background knowledge on Verus annotation generation; and (iv) a machine-parseable output format for the execution sequence.

%As discussed in \Cref{sec:overview}, dependencies among $G_i$'s permit a canonical order: $G_1 \rightarrow G_2 \rightarrow G_3 \rightarrow G_4$, but we can skip some unnecessary $G_i$'s. Thus, the planner agent instructs the LLM to determine whether each module should be executed or omitted, producing a subchain of this sequence as the final execution plan. 

The prompt is designed with several guiding principles:
\begin{enumerate}
    \item Invoke $M_1$ if the data structure can be represented by a sequence, set, or map.
    \item Invoke $M_2$ when non-trivial relationships exist among fields (e.g., range constraints or arithmetic relations).
    \item If $M_2$ is invoked, also execute $M_4$ to add type invariants into the proof context.
\end{enumerate}

%\paragraph{Planner}
 %The full details are presented in \Cref{sec:genfunc}.
%Given a to-be-verified data structure, the planner agent instructs the LLM to select and order the generation modules. Note that it is 

%Its prompt comprises four blocks: (i) a concise role description of the planner, (ii) a catalog of the available generation modules—one per annotation artifact discussed above, (iii) planning constraints that eliminate invalid orders by enforcing dependencies (e.g., generate the view and type invariant before method specifications; schedule proof blocks last), and (iv) a machine-parseable output format to ensure deterministic downstream execution. The full details are presented in \Cref{sec:approach}.

\noindent
For the ring buffer, the planner emits the full execution sequence, namely $\langle M_1, M_2, M_3, M_4\rangle$. 

%With this choice made, we next describe how \toolname{} synthesizes each component, beginning with the view generation; the remaining three components follow an analogous generation procedure.

\paragraph{Module Invocation} Following the generated plan, the framework sequentially invokes each module $M_i$ (Lines 6–10). The function \textbf{ExecWithSampling} takes as input: (1) a module to be executed; (2) the code to be verified; and (3) the unit test suite.  It executes the module, and returns the annotated code. To improve robustness, it asks for \texttt{n} samples during each module invocation~\cite{sampling}. It then selects the result yielding the maximum number of successfully verified functions.  We describe the view generation module in detail below, and then describe some differentiating aspects of the other modules, which are otherwise similar.

\paragraph{View Generation (and Refinement) Module}
To generate a view, our prompt (\Cref{fig:prompt-view-gen}) is structured as follows:
\begin{enumerate}
\item \emph{Objective.} This part articulates the task of generating a \texttt{View} implementation.%—to synthesize a minimal and semantically meaningful.
\item \emph{Verus Guidelines.}  This part provides conceptual foundations for generating a valid \texttt{View} implementation.
    As discussed in \Cref{sec:intro}, LLMs often struggle with Verus' syntax and verification-specific semantics. To mitigate this, our framework incorporates guidelines on Verus syntax and semantics imported from the Verus standard library documentation~\cite{verus-stdlib} and the official Verus tutorial~\cite{verus-tutorial}. For each different module in \toolname{}, a specific set of guidelines is added to the prompt. For \texttt{View} generation, for example, we include:
    \begin{itemize}
    \item An overview of the \texttt{View} trait (as introduced in \Cref{sec:background}).  
    \item Guidelines on how to manipulate Verus' logical types, which serve as foundational building blocks for providing a \texttt{View} implementation.  
    \item A chapter on Verus' specialized logical syntax.%, covering key constructs for writing specifications, such as using \texttt{==>} for implication.% and \texttt{&&&}/\texttt{|||} for conjunction/disjunction—to ensure syntactic validity in generated logical expressions.
    \end{itemize}

    \item \emph{Step-by-Step Instructions.} This part provides step-by-step procedural instructions on how to systematically generate a view.
    \item \emph{Examples}~\cite{incontextlearning}. This part provides examples of \texttt{View} implementations for other verified data structures (e.g., doubly-linked lists). These examples are kept the same across (and are disjoint from) all benchmarks.% and serve as reference points for few-shot in-context learning~\cite{incontextlearning}.  
    \item \emph{Code with Tests.} Finally, the prompt adds the code to be verified code along with the unit test suite.
    %This part gives examples of the view of other data-structure modules~\cite{incontextlearning}, we keep this part the same for all benchmarks.
\end{enumerate}


%Notebly, As point out in the challenge paragraph of \Cref{sec:intro}, LLM does not grasp the Verus syntax. To alliveate this isseue, we also add syntax guidance. The syntax guidance are built in chapters. We construct each chapter from the Verus standard library~\cite{?} and the Verus tutorial~\cite{?}. For example, we have syntax guidelines in manipulating Verus' axoiomic logical types~(\texttt{Seq<T>}, \texttt{Set<T>}, and \texttt{Map<K, V>}), we also have syntax guidelines in Verus' featured syntax in writing logical expressions~(e.g., use \texttt{==>} instead of the standard \texttt{->} for logical implication), these chapters will be added while generating the \texttt{View} trait, since the \texttt{View} trait requires building a view function which is some logical expression that built on top of in Verus' built-in logical types.


\begin{figure}[htbp]
\begin{lstlisting}[basicstyle=\fontsize{7pt}{7pt}\normalfont\ttfamily,frame=single,breaklines=true,escapechar=@,breakindent=0pt,xrightmargin=3em,xleftmargin=3em]
@\textbf{{[Objective]}}@
You are an expert in Verus (verifier for rust). Your task is to generate a View trait for the given data-structure module...
@\textbf{{[Verus Guidelines]}}@
Background on View and relevant Verus features...
@\textbf{{[Step-by-Step Instructions]}}@
1. Infer what should be contained in the View trait.
2. Generate the view based on the inferred information.
@\textbf{{[Examples]}}@
The View trait for other data structures, e.g., doubly-linked list.
@\textbf{{[Code and Tests]}}@
Code to be generated a View trait...
\end{lstlisting}
\caption{The Prompt for View Generation}
\label{fig:prompt-view-gen}
\end{figure}

\noindent
For the RingBuffer example, LLMs sometimes produce a trivial Cartesian-product view, which includes all of the data fields in RingBuffer (see \Cref{fig:rb-viewgen}), i.e., the view function simply maps the buffer content \texttt{self.ring} to its logical abstraction \texttt{self.ring.view()} and converts \texttt{head} and \texttt{tail} into logically defined natural numbers. 
This view function, while syntactically valid, fails to abstract away the circular structure of the ring buffer. As a result, the subsequent specification generation becomes significantly more complex: we have to explicitly reason about the arithmetic relationship between \texttt{head} and \texttt{tail} when specifying the behavior of operations such as \texttt{enqueue} and \texttt{dequeue}. %Consequently, this over-detailed view makes the specifications for these operations unnecessarily complicated and harder to verify.

\begin{figure}
\begin{lstlisting}[language=Rust,style=boxed,showstringspaces=false,basicstyle=\small\ttfamily]
impl View for RingBuffer<T> {
    type V = (Seq<T>, nat, nat);
    closed spec fn view(&self) -> Self::V {
        (self.ring.view(), head as nat, tail as nat) 
    }
}
\end{lstlisting}
\caption{Result of View Generation for RingBuffer}
\label{fig:rb-viewgen}
\end{figure}

To address this issue, we introduce an additional \emph{view refinement} step that prompts the LLM to reconsider the generated \texttt{View} and reconstruct it with fewer, more abstract components. This module encourages the model to focus on the conceptual essence of the data structure rather than on reproducing its concrete implementation. The refinement process substantially simplifies subsequent specification and proof generation tasks.

The \emph{view refinement} prompt follows the same structure as \Cref{fig:prompt-view-gen}, but emphasizes abstraction, minimality, and logical coherence over completeness. After applying this refinement, the model typically is able to produce an improved \texttt{View} implementation, like the one shown in \Cref{fig:ring-buffer-verified}.%, which succinctly captures the buffer’s logical content and capacity while abstracting away low-level representation details.

%Finally, we remark that the generation of the type invariant and specifications follows a similar pipeline; their detailed descriptions are presented in \Cref{sec:approach}.

\paragraph{Other Modules}
The prompts for the other modules share a similar structure to that shown in \Cref{fig:prompt-view-gen}. Each module prompt begins with a task description, followed by module-specific Verus guidelines.  These guidelines are specialized to the individual modules.

The \textit{Type Invariant Module} is focused on generating type invariants, which were introduced in \Cref{sec:background}.  In this module, the LLM is instructed to focus on common patterns such as range constraints, capacity checks, and arithmetic relationships between fields.%, not only we describe how type invariant works as in \Cref{sec:background}, but also instructs LLM to focus on common usages of type invariant, i.e., range constraints, capacity check, or field value arithematic bounds.

The \textit{Specification Module} is focused on generating specification functions, preconditions and postconditions.  We add mutability guidelines to ensure compatibility with Rust’s ownership and mutability system in \texttt{requires} and \texttt{ensures} clauses. For example, the LLM is prompted to use \texttt{old($\cdot$)} in \texttt{requires} to reference a variable's value prior to function invocation. %we will also add mutability guidelines, to tell LLM how to make Verus specifications compatible with Rust mutability type system in the \texttt{requires} and \texttt{ensures}. For instance, the LLM should write \texttt{old($\cdot$)} in \texttt{requires} to show that this refers to the variable before invoking this function.

The \textit{Proof Block Module} handles proof blocks and loop invariants.  The LLM is directed to construct the appropriate proof blocks and loop invariants, while explicitly incorporating type invariants into the proof context or applying relevant lemmas when necessary. %we will ask LLM do not forget to add type invariant into the context by calling \texttt{use\_type\_invariant}. 




%This section demonstrates the detail of the first stage, i.e., generating the annotations. As demonstrated in \Cref{sec:overview}, generating these annotations requires the generation of four different parts of annotations: ($P_1$) the \texttt{View} trait, ($P_2$) the type invariant, ($P_3$) the specifications (pre- and post-conditions), and ($P_4$) the proof blocks. Each part $P_i, i\in \{1,\dots,4\}$ corresponds to of a dedicated generation module $G_i$, in charge of generating all annotations in $P_i$. 



%Below, we demonstrate the details of the planner agent. As demonstrated in \Cref{sec:overview}, due to the depedencies between $P_i$'s, we can always arrange the generation order to be ($G_1$) $\rightarrow$ ($G_2$) $\rightarrow$ ($G_3$) $\rightarrow$ ($G_4$). Thus, the planner agent just looks at whether each step should be executed or emitted, and outputs a subchain of ($G_1$) $\rightarrow$ ($G_2$) $\rightarrow$ ($G_3$) $\rightarrow$ ($G_4$) as the exeuction order. The planner instructs LLM to decide whether to execute each step. The prompt is featured with several principles when a module is necessary, as below.
%Then, we invoke the generation modules in the plan one by one (Lines 6--10). Each module $G_i$ instructs the LLM to generate all annotations in $P_i$. Note that we have introduced $G_1$ (the generation of \texttt{View} trait) in \Cref{sec:overview}. The prompt in modules $G_{2,3,4}$ follows the same structure as \Cref{fig:prompt-view-gen}. All prompts of $G_{2,3,4}$ start with a task description, describing the corresponding part of annotations. Then follows the guidline information. Since we are writing annotaionts, the guidance of all $G_i$ consists of how to write a logical expression in Verus. Below, we list some additional features in the Verus guidline part.



% All $G_i$ modules also include general guidelines for writing logical expressions in Verus, since all annotations are formulated as logical constructs. Modules $G_{2,3,4}$ are treated as single-step modules and therefore omit step-by-step instructions. Each module also includes in-context learning examples to illustrate the desired annotation style and structure. In the end, we add the code with test suite into the prompt.
%Note that all $G_i$ also additionally consists of the guideline in writing logical expressions in Verus, because all annotations are logical expressions.
%Then, for the step-by-step instructions, all modules of $G_{2,3,4}$ are recognized as a one-step module, thus does not have this part. Finally, we add in-context learning examples for each module.

\section{Stage 2: Repairing the Annotations}
\label{sec:fixfunc}

Unfortunately, a single generation run rarely yields fully verified code. LLMs often struggle with Verus' syntax and rules. For example, the following is an \emph{incorrect} specification that was generated for the \texttt{enqueue} method of the \texttt{RingBuffer} example:
\begin{lstlisting}[language=Rust,style=boxed,showstringspaces=false,basicstyle=\small\ttfamily]
pub fn enqueue(&mut self, val: T) -> (succ: bool)
ensures
    succ == !old(self).is_full(),
    ... // omitted
{ /* omitted */ }
\end{lstlisting}
The intent is reasonable: \texttt{enqueue} succeeds if and only if the original buffer is not full. However, \texttt{is\_full} is an \emph{executable} function (see Line~46 of \Cref{fig:ring-buffer-verified}) and therefore cannot be called from a specification context.
The verifier reports:
\begin{lstlisting}[basicstyle=\fontsize{8pt}{7pt}\normalfont\ttfamily,frame=single,breaklines=true,escapechar=@,breakindent=0pt]
Error: cannot call the executable function is_full in annotation.
\end{lstlisting}

%\paragraph{Repair Loop}
To handle such issues, \toolname{} executes an {iterative repair loop}. In each iteration, it: ($i$) selects an error reported by the verifier; ($ii$) applies a corresponding predefined repair module; and ($iii$) reruns verification. The process continues until either all errors are eliminated or a preset iteration budget is exhausted. \toolname{} provides a suite of repair modules, each targeting a specific class of errors. Error messages are routed to the appropriate module via lightweight pattern matching over the verifier's error message.
%Specifically, to support verification of unit tests, \toolname{} integrates a special repair module that strengthens the method specifications when the verifier identifies an assertion failure in the input tests (details in \Cref{sec:fixfunc}).

For the error reported above, \toolname{} invokes a dedicated mode-repair module, which provides instructions for fixing this mode misuse~(Figure~\ref{fig:prompt-mode-repair}). After applying this module, a new specification is generated.  For example, one possible solution (one we actually observed) is to create a specification version of the \texttt{is\_full} function and use it to replace all occurrences of \texttt{is\_full}. With this change, the error no longer occurs, and the workflow can continue.
\begin{lstlisting}[language=Rust,style=boxed,showstringspaces=false,basicstyle=\small\ttfamily]
spec fn is_full_spec(&self) -> bool { self@.0.len() == self@().1 - 1 }
\end{lstlisting}



%For example, \toolname{} has a module for repairing mode errors illustrated above. The module designs a dedicated prompt for fixing these errors. [Summarize the prompt in this place]. 

\begin{figure}[htbp]
\begin{lstlisting}[basicstyle=\fontsize{7pt}{7pt}\normalfont\ttfamily,frame=single,breaklines=true,escapechar=@,breakindent=0pt,xrightmargin=3em,xleftmargin=3em]
@\textbf{{[Objective]}}@
Fix the error for the following code. The error indicates that an executable function is called from annotations or vice versa...
@\textbf{{[Relevant Background]}}@
Background on mode and relevant Verus Features
@\textbf{{[Code with Error, and Tests]}}@ ...
@\textbf{{[Instructions]}}@ Make one of these changes:
1. Adjust the function to be compatible with the calling context
...
\end{lstlisting}
\caption{The Prompt for Repairing the Misuse of Specification and Executable Functions}
\label{fig:prompt-mode-repair}
\end{figure}

\begin{small}
\begin{algorithm}[htbp]
%\SetAlgoLined
\SetAlgoNoLine 
\SetKwProg{Fn}{Function}{}{end}
\caption{The Repair Stage}
\label{fig:fixfunc}
\textbf{Parameter.} {\texttt{n}: Number of samples generated in each module invocation.}\\
\textbf{Parameter.} {\texttt{m}: Maximum number of iterations in the repair loop.}\\
\textbf{Given.} {\texttt{rep\_modules}: A pre-defined set of repair modules.}\\
\textbf{Input.} {\texttt{code}: the code after generating annotations; \texttt{test}: unit test suite}\\
\textbf{Output.} {\texttt{code} with annotations repaired}\\
\Fn{{{\rm \fixfunc{}}}{\rm (}{\tt code}, {\tt test}{\rm )}}{
\For{{\tt r} $\gets$ {\tt 1} \KwTo {\tt m}}{
    {\tt err} $\gets$ {\bf TryVerify}{\rm (}{\tt code}, {\tt test}{\rm )}\\
    \If{{\tt err is None}}{\Return {\tt Success, code}}{}
    {\tt cur\_module} $\gets$ {\tt default\_module}\\
    \ForEach{{\rm (}{\tt pattern}, {\tt module}{\rm )} $\in$ {\tt rep\_modules}}{
    \If{{\tt err} {\rm matches} {\tt pattern}}{
        {\tt cur\_module} $\gets$ {\tt module}\\
        \textbf{break}
    }{}
    }
    {\tt code} $\gets$ {\textbf{ExecWithSampling}}({\tt cur\_module}, {\tt code}, {\tt test}, {\tt n})
}
\Return {\tt Failure, code}
}
\end{algorithm}
\end{small}

\Cref{fig:fixfunc} shows the repair algorithm in detail.  \toolname{} provides a predefined set of repair modules, denoted as \texttt{rep\_modules}. Each module is associated with a regular expression \texttt{pattern} specifying the class of error messages that it addresses. In the current implementation, \toolname{} includes repair modules for: (1) misuse between specification and executable functions (as demonstrated above), (2) inconsistencies with Rust’s mutability type system, (3) violations of preconditions and postconditions, (4) arithmetic overflow or underflow, (5) mismatches between logical and Rust types, (6) test assertion failures, and (7) a fallback module, \texttt{default\_module}, which prompts the LLM to attempt a repair based directly on the provided error message.

Note that the repair module for test assertion failures requires the LLM to do interprocedural analysis, something largely absent from previous work.  The repair module inspects the failing assertion, identifies the method invoked immediately before it, and then instructs the LLM to strengthen that method's postcondition so as to ensure the assertion holds\ifarxiv{} (Table~\ref{tab:repair-heuristics})\fi.

%\livia{Shan commented that this interprocedual generation should be highlighted
%}

The repair process proceeds iteratively, with a maximum number of iterations defined by the parameter \texttt{m}. In each iteration, the system first attempts to verify the current implementation via \textbf{TryVerify} (Line 8), which takes the code and test suite as input and returns the highest priority error message \texttt{err} produced by Verus. If no error is reported, verification succeeds and the code is returned (Lines 9--10). Otherwise, \texttt{err} is matched against all module patterns; if a match is found, the corresponding repair module is executed (Lines 13--15). If no match is found, the \texttt{default\_module} is invoked. Each module execution includes {\tt n} sampling attempts to improve robustness~(Line 16). %In addition to syntax-level fixes, some repair modules operate semantically—for instance, the \emph{repair test assertion failure} heuristic inspects the test assertion that triggered an error, locates the method invoked immediately before the failing assertion, and strengthens that method's postcondition so that the assertion can hold. 
If the implementation remains unverified after \texttt{m} iterations, the system reports failure and returns the final unverified result~(Line 17).

%(1) Rust's mutability type system, (2) misuse between specification and executable functions, (3) pre- and post-condition failures, (4) arithmetic overflow, (5) misuse between logical types and Rust types, and a last resort module \texttt{default\_module} which simply ask LLM to repair the error based on the provided message.

%The repair stage is performed in iterations, where the maximum number of iterations is a pre-defined parameter {\tt m}. In each iteration, we first try to verify the current implementation {\tt code} (Line 8). The function {\bf TryVerify} takes the code and the test suite as the input, and outputs the \emph{first} error message {\tt err} reported by Verus. If there is no error message, which means the code has been verified, we return the {\tt code} (Lines 9--10). Otherwise, we match the first error {\tt err} with the patterns of all repairing modules, if any of them mathches, then execute the corresponding module (Lines 13--15), if neither matches, perform the default repair  \texttt{default\_module}. The execution of each module is also coupled with sampling {\tt n} times.

%If the code is still unverified after {\tt m} rounds, report failure.